\chapter{Building CNA}
\label{ch:building-cna}

\section{Prerequisites}

CNA targets C\texttt{++}23 and CMake 3.20+. On Linux, the minimum toolchain is GCC~12+ or
Clang~15+; two sibling directories must be available to CMake alongside the CNA checkout
itself: \texttt{../sharp-runtime} (always required) and \texttt{../easy-gl} (only required
if you build the \texttt{EASYGL} backend). SDL3, SDL3\_image, and SDL3\_mixer are built from
vendored Git submodules by default, so no system SDL packages are required on a first build.

On Windows, any of MSVC~2022 (v17.8+), clang-cl, or MinGW-w64 (native or cross-compiled from
Linux) works; only \texttt{../sharp-runtime} is needed as a sibling dependency, and the same
vendored-submodule SDL story applies.

\subsection{Initializing submodules}

Before the first build, CNA's vendored SDL3/SDL3\_image/SDL3\_mixer submodules must be
populated:

\begin{lstlisting}[style=cnashell]
git submodule update --init --recursive
\end{lstlisting}

A GitHub ``Download ZIP'' archive does \emph{not} include submodule contents. Building from
such an archive leaves \texttt{third\_party/SDL} empty, and CMake aborts at configure time
with a clear, specific error pointing at \texttt{cmake/ThirdPartySDL.cmake} rather than
failing confusingly deep in a missing-header compile error. The two ways around this are
cloning with Git and running the command above, or configuring with
\texttt{-DCNA\_USE\_SYSTEM\_SDL=ON} to use system-installed SDL3 packages instead of the
vendored submodules.

\section{The default build: Linux, EasyGL backend}

\begin{lstlisting}[style=cnashell]
git submodule update --init --recursive
cmake -S . -B build -DCNA_GRAPHICS_BACKEND=EASYGL
cmake --build build --target CNA CnaTests
ctest --test-dir build --output-on-failure
\end{lstlisting}

\texttt{EASYGL} is the most mature backend by the project's own account (Chapter~\ref{ch:easygl-backend}),
so it is the natural default for a first build. Swapping to \texttt{SDL\_RENDERER} is a
one-flag change:

\begin{lstlisting}[style=cnashell]
cmake -S . -B build-sdlrenderer -DCNA_GRAPHICS_BACKEND=SDL_RENDERER
cmake --build build-sdlrenderer --target CNA CnaTests
\end{lstlisting}

This one-flag swap is the whole point of the backend abstraction covered in
Chapter~\ref{ch:backend-architecture} --- the same \texttt{CNA} / \texttt{CnaTests} build
targets exist regardless of which of the fourteen backends in
Table~\ref{tab:backend-option-list} you select.

\begin{table}[h]
\centering
\begin{tabular}{ll}
\toprule
\textbf{\texttt{CNA\_GRAPHICS\_BACKEND} value} & \textbf{Platform constraint} \\
\midrule
\texttt{SDL\_RENDERER} & Cross-platform \\
\texttt{EASYGL}        & Cross-platform (default focus) \\
\texttt{BGFX}          & Cross-platform \\
\texttt{VULKAN}        & Cross-platform (Vulkan-capable GPU/driver) \\
\texttt{WEBGPU}        & Cross-platform, experimental \\
\texttt{HEADLESS}      & Cross-platform, no window/present \\
\texttt{SOFTWARE}      & Cross-platform, CPU rasterizer \\
\texttt{D3D11}         & Windows-only \\
\texttt{D3D12}         & Windows-only \\
\texttt{CANVAS}        & Emscripten-only \\
\texttt{ASCII}         & Cross-platform, SDL-windowed glyph-grid decorator \\
\texttt{DX3}           & Cross-platform (via \texttt{../free-direct}) \\
\texttt{SDL\_GPU}      & Cross-platform (Vulkan driver on Linux; Chapter~\ref{ch:sdlgpu-backend}) \\
\bottomrule
\end{tabular}
\caption{Every value accepted by \texttt{-DCNA\_GRAPHICS\_BACKEND}. \texttt{D3D9} is also
Windows-only via cross-compilation; see Chapter~\ref{ch:direct3d-backends} for its own build
invocation, which additionally requires \texttt{-DCNA\_BUILD\_TESTS=ON}.}
\label{tab:backend-option-list}
\end{table}

\section{The full CMake option reference}
\label{sec:cmake-option-reference}

\texttt{-DCNA\_GRAPHICS\_BACKEND} is the option most build invocations in this book set
explicitly, but CNA's top-level \texttt{CMakeLists.txt} declares several more, every one of
them read directly from that file rather than assumed from memory:

\begin{center}
\begin{tabularx}{\textwidth}{llX}
\toprule
\textbf{Option} & \textbf{Default} & \textbf{Effect} \\
\midrule
\texttt{CNA\_USE\_CCACHE} & \texttt{ON} & Auto-detects \texttt{ccache} and wires it in as the
  C/C\texttt{++} compiler launcher for both CNA and \texttt{sharp-runtime}, unless a launcher
  is already set. \\
\texttt{CNA\_BUILD\_TESTS} & \texttt{ON} & Builds the \texttt{CnaTests} GoogleTest target and
  registers CTest entries. \\
\texttt{CNA\_NOXNA} & \texttt{OFF} & Enables an opt-in, project-specific extended render
  pipeline (see below) --- unrelated to the \texttt{NOXNA} marker macro from
  Chapter~\ref{ch:design-philosophy}. \\
\texttt{CNA\_DEVICES} & \texttt{OFF} & Enables CNA-specific device/sensor extensions beyond
  XNA 4.0 (battery, camera, clipboard, and similar). \\
\texttt{CNA\_ENABLE\_NET} & \texttt{ON} & Builds \cnans{Net} and \cnans{GamerServices}
  networking; pulls in the vendored ENet dependency when on. \\
\texttt{CNA\_USE\_SYSTEM\_SDL} & \texttt{OFF} & Uses \texttt{find\_package} for system SDL3
  packages instead of the vendored submodules (see above). \\
\texttt{CNA\_TEST\_DISPLAY} & \texttt{":0"} & The \texttt{DISPLAY} value baked into every
  GPU/window-creating CTest's environment; override to point the suite at a virtual display
  (e.g.\ Xvfb) instead of a real desktop. \\
\texttt{CNA\_SANITIZE} & \emph{(empty)} & Comma-separated sanitizer list (e.g.\
  \texttt{address,undefined}) added to both compile and link flags project-wide; diagnostics
  only, never for a shipped build. \\
\bottomrule
\end{tabularx}
\end{center}

Every one of these is independent of \texttt{-DCNA\_GRAPHICS\_BACKEND}
(Table~\ref{tab:backend-option-list}).

\begin{sourcenote}
\textbf{Two same-shaped names, two unrelated things.} \texttt{CNA\_NOXNA} (a CMake
\texttt{option()}, off by default) and \texttt{NOXNA} (the marker macro from
\S\ref{sec:noxna}) share a name fragment purely by coincidence of naming convention, not by
design --- confusing the two is an easy real mistake to make when skimming the source. The
marker macro tags an individual declaration as ``not part of stock XNA 4.0'' and is present in
every build regardless of any flag. \texttt{CNA\_NOXNA} the CMake option instead gates an
entirely separate, genuinely optional feature: a settings-bag-only ``extended render
pipeline'' (\cnaclass{RenderPipelineSettings}, under \cnans{CNA::Graphics}, guarded behind
\texttt{\#ifdef CNA\_NOXNA} in its own header) covering HDR exposure/gamma, tonemapping mode,
shadow quality, and PBR material settings. Its own header comment is explicit about scope:
it ``does not perform any rendering itself'' --- it is configuration a backend can read and
act on, with backend support for each individual feature not yet implemented across the board
as of the snapshot this book draws from. Building with plain \texttt{cmake -S . -B build}
(the option's own \texttt{OFF} default) never compiles this pipeline in at all.
\end{sourcenote}

\section{Windows, natively and cross-compiled}

On real Windows, \texttt{SDL\_RENDERER} is selected automatically when no backend is
specified, and SDL is built from the vendored submodule with no pre-built binaries or
\texttt{CMAKE\_PREFIX\_PATH} configuration needed:

\begin{lstlisting}[style=cnashell]
git submodule update --init --recursive
cmake -S . -B build-win -DCNA_GRAPHICS_BACKEND=SDL_RENDERER
cmake --build build-win --target CNA CnaTests
\end{lstlisting}

The same target can be reached without a Windows machine, by cross-compiling from Linux with
MinGW-w64 --- the toolchain every Windows-targeting CTest binary in this book's Part~IV
chapters (D3D9/D3D11/D3D12, and the plain \texttt{SDL\_RENDERER} Windows target) is actually
verified through:

\begin{lstlisting}[style=cnashell]
sudo apt install mingw-w64

git submodule update --init --recursive
cmake -S . -B build-windows \
      -DCMAKE_TOOLCHAIN_FILE=cmake/toolchains/mingw-w64.cmake \
      -DCNA_GRAPHICS_BACKEND=SDL_RENDERER
cmake --build build-windows --target CNA CnaTests
\end{lstlisting}

Chapter~\ref{ch:direct3d-backends} and this book's cross-platform verification chapter
cover what happens after the cross-compiled \texttt{.exe} exists: running it under Wine, with
DXVK (D3D9/D3D11) or vkd3d-proton (D3D12) translating the Direct3D calls to Vulkan on a real
GPU, which is how this project verifies Windows-only backends without ever touching a
Windows machine.

\section{Building for Android and the Web}

Neither of these two targets gets its own \texttt{cmake/toolchains/} file the way
\texttt{mingw-w64.cmake} does for Windows --- both instead rely on toolchain infrastructure
that ships with the NDK and the Emscripten SDK respectively, invoked directly.

\subsection{Android (NDK)}

Grounded in the project's own \texttt{docs/devices-build.md}, a real cross-compile invocation
looks like this:

\begin{lstlisting}[style=cnashell]
cmake -S . -B cmake-build-android -G Ninja \
  -DCMAKE_TOOLCHAIN_FILE="$HOME/Android/Sdk/ndk/30.0.14904198/build/cmake/android.toolchain.cmake" \
  -DANDROID_ABI=arm64-v8a \
  -DANDROID_PLATFORM=android-24 \
  -DCNA_BUILD_TESTS=OFF

cmake --build cmake-build-android --target CNA -j"$(nproc)"
\end{lstlisting}

Three details worth knowing before running this yourself, each one a real, documented finding
rather than a generic Android-CMake tip: the exact NDK version path
(\texttt{30.0.14904198}) is whatever happens to be installed in a given environment, not a
fixed requirement --- check \texttt{ls~\textasciitilde{}/Android/Sdk/ndk/} first rather than
assuming either that it exists or that this exact version is what you have.
\texttt{-DCNA\_BUILD\_TESTS=OFF} is not optional on this path today: \texttt{googletest} was
never configured for the NDK toolchain, so only the \texttt{CNA} static library itself
cross-compiles, not \texttt{CnaTests}. And this specific invocation is a
\textbf{compile-only} verification --- no APK packaging and no emulator/device run happen as
part of it; Chapter~\ref{ch:android-ndk} covers the further steps (the NDK's own
\texttt{llvm-nm}, not the host's plain \texttt{nm}, is what actually confirms
platform-conditional code compiled in) that turn a successful compile into a running,
on-device verification.

\subsection{Web (Emscripten)}

The \texttt{CANVAS} backend (Chapter~\ref{ch:canvas-ascii-backends}) is Emscripten-only, hard-gated
at configure time the same way the Direct3D backends are hard-gated to Windows. Building it
follows the standard Emscripten CMake pattern --- \texttt{emcmake} wraps the configure step,
\texttt{emmake} (or a plain \texttt{cmake~--build}, once configured) wraps the build --- applied
to the same \texttt{-DCNA\_GRAPHICS\_BACKEND} flag every other backend uses:

\begin{lstlisting}[style=cnashell]
emcmake cmake -S . -B build-web -DCNA_GRAPHICS_BACKEND=CANVAS
cmake --build build-web
node build-web/CnaTests.js
\end{lstlisting}

This is a real, verified path, not a hopeful sketch: the project's own tracking records a
genuine \texttt{emcmake} / \texttt{emcc}~6.0.2 configure-and-build succeeding, with the resulting
\texttt{CnaTests.js} linking cleanly and a backend-agnostic GTest suite (35 of 35
\cnaclass{Rectangle}-family tests) genuinely passing under \texttt{node}. Getting there
surfaced a real, pre-existing bug outside CNA's own tree entirely: \texttt{sharp-runtime}'s
\texttt{FileSystemWatcher.hpp} declared three \texttt{inotify}-specific fields
unconditionally, even though they were only ever \emph{used} under \texttt{\#ifdef
\_\_linux\_\_} --- a guard \texttt{emcc} does not define, so the fields existed but nothing
initialized them, a genuine latent bug independent of Emscripten that this cross-compile
attempt is what actually surfaced. Fixed directly in \texttt{sharp-runtime}, project-owner
approved. What this build does \emph{not} yet prove is anything about a real browser: this
project's own dev loop has no DOM and no \texttt{CanvasRenderingContext2D} at all, so
\texttt{node}-based verification and real-browser verification remain two genuinely different
claims, covered in full in Chapter~\ref{ch:canvas-ascii-backends}.

\section{System-installed SDL, if you prefer it}

Rather than the vendored submodules, system SDL3/SDL3\_image/SDL3\_mixer packages can be
used instead:

\begin{lstlisting}[style=cnashell]
cmake -S . -B build -DCNA_USE_SYSTEM_SDL=ON -DCNA_GRAPHICS_BACKEND=SDL_RENDERER
\end{lstlisting}

This calls \texttt{find\_package(SDL3 REQUIRED)},
\texttt{find\_package(SDL3\allowbreak{}\_image REQUIRED)}, and
\texttt{find\_package(SDL3\allowbreak{}\_mixer REQUIRED)}, and requires those development
packages to already be installed system-wide.

\section{Linux system dependencies beyond SDL}

Two optional native libraries extend specific subsystems and are not vendored:

\begin{itemize}
  \item \textbf{FFmpeg} (\texttt{libavcodec-dev libavformat-dev libavutil-dev
        libswresample-dev}) is required for \cnaclass{VideoPlayer} (video decoding).
        \texttt{libswscale-dev} is notably \emph{not} required even though FFmpeg's own video
        pipeline usually pairs with it --- CNA implements YUV$\to$RGBA conversion internally,
        so only the runtime \texttt{libswscale8} library needs to be present, not its headers.
  \item \textbf{Draco} (\texttt{libdraco-dev}) is optional and enables
        \texttt{KHR\_draco\_mesh\_compression} decoding in \texttt{GltfImportCore}. It is
        detected via CMake's \texttt{find\_package(draco CONFIG)}; when absent, importing a
        Draco-compressed glTF primitive throws a clear ``not supported'' error at import time
        rather than failing to build at all. Unlike \texttt{cgltf.h} and \texttt{stb\_image.h}
        (single-header vendored dependencies), Draco is a real multi-file library and is
        deliberately not vendored.
\end{itemize}

\section{Verifying a build}

CNA intentionally does not ship a bundled playable game demo as its primary verification
artifact --- the project prioritizes framework/runtime development over a shipping demo
executable. Two commands stand in for ``does this build actually work'':

\begin{lstlisting}[style=cnashell]
ctest --test-dir build --output-on-failure
cmake --build build --target hello-triangle-sdl
\end{lstlisting}

\texttt{ctest} runs the GoogleTest-based \texttt{CnaTests} suite (thousands of unit tests,
plus, on backends where a display is available, GPU pixel tests); \texttt{hello-triangle-sdl}
is the minimal environment/rendering-path smoke test mentioned throughout the backend chapters
in Part~IV.

\begin{sourcenote}
\textbf{On CI coverage.} As documented in CNA's own README, project CI is Linux-only and
partial: at the time this book's source material was gathered, CI ran only the
\texttt{Input} and \texttt{Devices} / \texttt{Sensors} GoogleTest suites --- not the full
roughly 4\,370-test unit suite, and not the roughly 490-test GPU pixel-test suite across the
four established graphics backends. Everything this book states about Graphics-layer
correctness is verified by running the suites locally and by hand, not by an automated
Graphics CI gate, and there is no Windows, macOS, or Android CI at all --- those platforms
are verified manually. Treat a green CI badge, if you see one, as covering a narrower slice
than ``the whole project passes.''
\end{sourcenote}
